
\element{partie1}{
    \begin{question}{verin-io}
        \evaluationProf[2]\\
        A partir de la description précédente, de la vidéo et de la documentation du vérin roto-linéaire fournie, proposer une liste des entrées et sorties nécessaires pour piloter un vérin de blocage.\\
        Compléter la section GVL avec vos réponse. Donner des adresses compatibles avec le type utilisé. \\
    \end{question}
}

\element{partie1}{
    \begin{question}{verin-missions}
        \evaluationProf[2]\\
        A partir de la vidéo et de la description du système, proposer une liste des missions que doit pouvoir réaliser un vérin de blocage roto-linéaire. Vous répondrez en déclarant le type de données \lstinline{TMissionVerinBlocage} dans votre programme.
    \end{question}
}

\element{partie1}{
    \begin{question}{verin-etats}
        \evaluationProf[2]\\
        Proposer la déclaration d'un type de données adaptées pour décrire les états possibles d'un vérin de blocage.\\
    \end{question}
}

\element{partie1}{
    \begin{question}{verin-declaration}
        \evaluationProf[2]\\
        A partir des questions précédentes, déclarer la classe \lstinline{CVerinBlocage} et compléter son contexte.\\
        \textit{Votre réponse contiendra la déclaration de la classe, les variables d'entrées, de sorties ainsi que les variables représentant les états et les missions.}\\
    \end{question}
}

\element{partie1}{
    \begin{question}{getSet}
        \evaluationProf[2]\\
        Donner l'implémentation des propriétés permettant d'accéder aux missions et états de la classe.\\
        \textit{Vous indiquerez également celles qui vous semblent inutiles en justifiant votre réponse.}
    \end{question}
}

\element{partie1}{
    \begin{question}{verin-corps}
        \evaluationProf[2]\\
        Proposer le contenu de la méthode permettant l'exécution de la mission reçue. (\textit{MExecutionMission})\\
    \end{question}
}

\element{partie1}{
    \begin{question}{setoutput}
        \evaluationProf[2]\\
        Proposer le corps de la fonction permettant de mettre à jour les sorties de la station d'alimentation.
    \end{question}
}

\element{partie1}{
    \begin{question}{verin-appel}
        \evaluationProf[2]\\
        Appeler un bloc fonctionnel \lstinline{oStation} de type \lstinline{CStationAlimentation} dans le programme principal.
    \end{question}
}



\pagebreak
\element{partie1}{
    \begin{question}{verin-missions-2}
        Proposer un grafcet de la mission permettant le blocage d'un volet.\\
        \evaluationProfGrille{2}{14}
    \end{question}
}
